%=============================================================================
% Section 9: Implementation Considerations
%=============================================================================
\section{Implementation Considerations}

\subsection{Runtime Interface}

The execution runtime exposes a minimal interface:

\begin{lstlisting}[style=pythonstyle]
interface ExecutionRuntime {
    # Execute a capability
    Result execute(CapabilityRequest request)

    # Query capability availability
    List<Capability> listCapabilities()

    # Health check
    HealthStatus healthCheck()
}
\end{lstlisting}

\subsection{Capability Request Format}

Capability requests follow a standardized format:

\begin{lstlisting}[style=jsonstyle]
{
    "capability": "generate_text",
    "version": "1.0",
    "input": {
        "prompt": "Explain quantum computing",
        "max_tokens": 500
    },
    "context": {
        "session_id": "abc123",
        "request_id": "req456"
    },
    "provider": {
        "name": "provider_a",
        "endpoint": "https://api.provider-a.com"
    }
}
\end{lstlisting}

\subsection{Result Format}

Results are normalized to a standard structure with execution metrics:

\begin{lstlisting}[style=jsonstyle]
{
    "capability": "generate_text",
    "status": "success",
    "output": {
        "text": "Quantum computing is...",
        "tokens_used": 450
    },
    "metadata": {
        "provider": "provider_a",
        "latency_ms": 1250,
        "translation_latency_ms": 15,
        "retry_count": 0,
        "timestamp": "2026-03-15T10:30:00Z"
    }
}
\end{lstlisting}

The \texttt{retry\_count} and \texttt{translation\_latency\_ms} fields provide visibility into execution behavior, enabling the Contact Layer to make data-driven routing decisions.
